\section{Related Work}
\textbf{Foothold planning.} Classical foothold selection scores candidate contact points
against local slope, roughness, and clearance heuristics computed from an elevation
map~\cite{singh2018survey}. Learned variants replace hand-tuned scoring with implicit terrain
encoders trained on simulated rollouts~\cite{mbeki2023legged}, improving generalization but
inheriting whatever noise is present in the input map.

\textbf{Whole-body and reactive control.} Whole-body controllers couple foothold choice to
torso stabilization through impedance or MPC formulations~\cite{nakashima2019whole,
delacroix2021mpc}, and reactive controllers detect stumbles via force or IMU anomalies and
trigger corrective swing trajectories~\cite{oduya2020reactive}. These methods are complementary
to ours: they act during or after a bad contact, whereas \textsc{ContactRefine} acts before
touchdown.

\textbf{Learning from experience.} Curriculum and domain-randomized reinforcement learning
have produced gaits that adapt to terrain roughness in simulation~\cite{papazian2022rl,
fontaine2021domain}, with sim-to-real transfer remaining the central bottleneck
\cite{halvorsen2022survey}. Our approach sidesteps large-scale simulation by learning directly
from real contact outcomes recorded on the target platform, at the cost of requiring an initial
teleoperated data collection phase.
